% Appendix: Evaluation Detail
% Moved from body to appendix to meet 15-page limit.

\section{Evaluation Detail}\label{app:eval-detail}

\subsection{Defect Discovery Register}\label{app:bug-cost}

Table~\ref{tab:bug-cost} catalogues every defect discovered during testing, organised by the tier at which it surfaced, together with the measured fix time and the assessed consequence had it reached flight day undetected.

\begin{table}[ht]
\centering
\caption{Defects discovered at each testing tier, with measured resolution time and impact if undetected. All fix times are from git commit timestamps; ``cost if missed'' is the assessed flight-day consequence.}
\label{tab:bug-cost}
\small
\begin{tabular}{@{}p{4.2cm}clp{4.5cm}@{}}
\toprule
\textbf{Defect} & \textbf{Tier} & \textbf{Fix Time} & \textbf{Cost if Missed} \\
\midrule
Geofence sign inverted (\texttt{pointPolygonTest} +ve = inside) & 1 & 15\,min & SSSI incursion; CAA regulatory violation \\
\addlinespace
Repulsive force reversed at NFZ boundary & 1 & 20\,min & Drone attracted \emph{into} NFZ \\
\addlinespace
Barometric drift $\to$ infinite landing loop & 1 & 30\,min & Battery drain; hover until failsafe \\
\addlinespace
Duplicate target queued (consecutive frames) & 1 & 10\,min & Repeated descent on same target \\
\addlinespace
CENTERING state missing timeout guard & 1 & 15\,min & State machine stuck indefinitely \\
\addlinespace
BGR/RGB colour inversion (IMX296 sensor) & 2 & 45\,min & Confidence drops to ${\sim}0.3$; mission failure \\
\addlinespace
FOV miscalibration ($f{=}7.0 \to 5.46$\,mm) & 2 & 60\,min & GPS estimates \SI{28}{\percent} off; landing outside \SI{10}{\metre} R07 radius \\
\addlinespace
Python 3.13 pyserial regression & 2 & 120\,min & Zero Pi--Cube communication \\
\addlinespace
\texttt{tflite-runtime} absent on 3.13 & 2 & 30\,min & Zero on-device inference \\
\addlinespace
Camera resource conflict (picamera2) & 2 & 20\,min & Stream or detection crash on startup \\
\addlinespace
\texttt{self.investigating} uninitialised & 2 & 5\,min & Python crash on first detection \\
\addlinespace
TFLite bbox pixel vs.\ normalised coords & 2 & 25\,min & Green boxes off-screen; GPS estimate \SI{100}{\metre}+ error \\
\addlinespace
GPS timing lag (100--200\,ms) & 3$^*$ & --- & 0.8--1.6\,m systematic offset at \SI{8}{\metre\per\second} \\
\midrule
\multicolumn{2}{@{}l}{\textbf{Total (Tiers 1--2)}} & \textbf{6\,hr 35\,min} & \\
\bottomrule
\multicolumn{4}{@{}l}{\footnotesize $^*$Discovered via DJI video analysis (Tier~3 proxy); mitigated by inverse-variance averaging and attitude compensation (\texttt{-{}-compensate-tilt}).} \\
\end{tabular}
\end{table}

\subsection{What Would Change in a Second Iteration}\label{app:what-would-change}

Five changes would yield the largest improvement if the project were repeated:

\begin{enumerate}
    \item \textbf{Start with Pi hardware from week~3, not week~16.} Seven of the thirteen integration defects (Table~\ref{tab:bug-cost}) were hardware-specific (Tier~2) and consumed \SI{4.9}{\hour} of compressed bench-day time. Earlier hardware access would have distributed these faults across 13 weeks instead of one day. The delay was partly a procurement issue (Pi~5 delivery) and partly a prioritisation error: the team assumed simulation maturity equated to hardware readiness. It did not---sensor calibration errors are invisible in simulation (Section~\ref{sec:sim-lessons}).

    \item \textbf{Run simulation and hardware development in parallel, not sequentially.} The simulation-first strategy consumed ${\sim}$10 of 16 weeks before hardware was powered on. While this produced a mature codebase, it compressed all integration testing into a single bench day---creating the exact time pressure the progressive framework was designed to avoid. A better approach: start simulation in week~1 \emph{and} power on Pi hardware in week~3. Run Tiers~1--3 concurrently. This requires accepting that the simulation code will be less polished, but it eliminates the ``integration cliff'' that the project encountered.

    \item \textbf{Automate regression tests in CI.} All 74 test scripts are run manually. A GitHub Actions pipeline running dry-run validation, TFLite model loading, geofence sign checks, and the 127 unit tests on every push would have caught at least 2 of the 5 Tier-1 defects (geofence sign inversion, duplicate target queue) before they reached the bench. The test suite's value is diminished by the fact that tests are only run on demand; continuous integration would convert the 74 scripts from documentation into active safety nets.

    \item \textbf{Record flight video in week~1 for training data.} The DJI video replay pipeline---built post-cancellation in a single session---produced the 4 key quantitative figures in this evaluation. Recording real aerial footage during \emph{any} early manual flight would have provided labelled training data weeks earlier, enabling a statistically meaningful held-out test set ($n \geq 50$) and eliminating the train/validation overlap concern (D5). This is the single most impactful missed opportunity: a 15-minute manual DJI flight in week~1 would have provided ground-truth data that took 16 weeks to obtain indirectly.

    \item \textbf{Deploy NCNN backend and train a multi-class detector from day~1.} NCNN benchmarks show a $4.5\times$ speedup over TFLite (\SI{72}{\milli\second} vs \SI{327}{\milli\second} full pipeline); the backend is implemented but has not yet been validated in the field. Pairing it with a multi-class detector (person/dummy/debris) would reduce operator verification load at the VERIFY state, addressing both D2 and D6.
\end{enumerate}

\subsection{Lessons Learned}\label{app:lessons-learned}

Four technical lessons emerged that generalise beyond this project:

\textbf{1. Invest in data collection infrastructure before model training.}
The synthetic dataset generator produced 300 images in ${\sim}\SI{2}{\minute}$, but the resulting model's mAP$_{50}{=}0.995$ is inflated by train/validation overlap (D5): only 16 of 366 training images (\SI{4.4}{\percent}) originate from real flight footage. The labelling tool (\texttt{label\_tool.py}) was built in week~16; building it in week~1 and collecting real frames during any early manual flight would have yielded a proper held-out test set ($n \geq 50$) and eliminated the need to treat mAP as an upper bound.

\textbf{2. Characterise the sensor pipeline end-to-end before writing application code.}
Three of the thirteen defects---BGR inversion (\SI{45}{\minute}), FOV miscalibration (\SI{60}{\minute}), and GPS timing lag (mitigated but not fully eliminated)---were sensor-pipeline issues that propagated silently into higher-level logic. A dedicated ``sensor characterisation sprint'' in week~1---measuring colour encoding, FOV at multiple distances, GPS latency under motion, and lens distortion---would have prevented these faults and saved ${\sim}\SI{2}{\hour}$ of bench-day debugging.

\textbf{3. Simulation fidelity has diminishing returns; real-world data does not.}
SITL simulation caught 5 of 13 defects (Tier~1) but assumes perfect GPS, zero wind, and synthetic camera views. The DJI video replay pipeline---built in one session post-cancellation---produced 4 quantitative figures (Figures~\ref{fig:conf_vs_alt}--\ref{fig:perf_analysis}) and the CEP$_{50}{=}\SI{2.3}{\metre}$ GPS characterisation that simulation could never yield. Future projects should prioritise replaying real sensor data through the pipeline over building higher-fidelity simulators.

\textbf{4. Budget for weather and schedule one flight day per week, not per project.}
A single scheduled flight day creates a binary outcome: either the weather cooperates and the project obtains flight data, or it does not. Scheduling weekly flight slots---even short 30-minute windows---would have provided multiple opportunities and allowed the progressive test tiers to be executed across several sessions rather than compressed into a single day.

\textbf{5. Test quantity does not equal test quality.}
The project produced 74 test scripts totalling ${\sim}$5\,500 LOC. Of these, 12 caught actual defects. The remaining 62 scripts served documentation, calibration, or development purposes---valuable, but not defect-catching. In a second iteration, the team should distinguish between \emph{validation tests} (expected to catch defects) and \emph{characterisation tests} (expected to produce data), and prioritise the former. A CI pipeline running even 10 critical validation tests on every commit would provide more safety assurance than 74 scripts run manually and sporadically.

\textbf{6. The ``sim-to-real gap'' is really a ``sim-to-real cliff'' without iterative hardware access.}
The project's multi-fidelity simulation framework was designed to narrow the sim-to-real gap incrementally across five tiers. In practice, Tiers~1--3 were executed over 10 weeks, and Tiers~4--5 were compressed into a single day that was then cancelled by weather. The progressive framework \emph{as designed} is sound, but its value depends on continuous access to real hardware throughout the project, not just at the end. Without iterative hardware testing, the simulation framework provides a false sense of readiness that collapses on first contact with physical reality---exactly as the 7 hardware-specific defects demonstrated.
